\section*{NeurIPS Paper Checklist}

\begin{enumerate}

\item {\bf Claims}
    \item[] Question: Do the main claims made in the abstract and introduction accurately reflect the paper's contributions and scope?
    \item[] Answer: \answerYes{}
    \item[] Justification: The abstract and Section 1 state the scope (a pilot on a 40-instance slice per substrate, one sweep per arm), the four contributions, and that all measurements are exploratory; Section 7 lists the limitations that bound them.

\item {\bf Limitations}
    \item[] Question: Does the paper discuss the limitations of the work performed by the authors?
    \item[] Answer: \answerYes{}
    \item[] Justification: Section 7 covers slice size, single sweeps, the adapter confound of the cross-stack comparison, the post hoc gated arms and their selection leakage, trajectory length relative to public scaffolds, the weak ungated baseline, and event counts as lower bounds; Section 5.4 discloses the grading correction made after unblinding.

\item {\bf Theory assumptions and proofs}
    \item[] Question: For each theoretical result, does the paper provide the full set of assumptions and a complete (and correct) proof?
    \item[] Answer: \answerNA{}
    \item[] Justification: The paper contains no theoretical results.

\item {\bf Experimental result reproducibility}
    \item[] Question: Does the paper fully disclose all the information needed to reproduce the main experimental results of the paper to the extent that it affects the main claims and/or conclusions of the paper (regardless of whether the code and data are provided or not)?
    \item[] Answer: \answerYes{}
    \item[] Justification: Section 3 specifies the instrument (commit per mutating call, replay of previously-passing tests in the pinned container, attribution, validation gates), Section 4 the substrates, slice rules, harness, tools, model identifiers, cost caps, outcomes and pre-declaration, and Appendix B the analysis plan; the defect catalogue (Appendix A) records the measurement pitfalls a replication must avoid.

\item {\bf Open access to data and code}
    \item[] Question: Does the paper provide open access to the data and code, with sufficient instructions to faithfully reproduce the main experimental results, as described in supplemental material?
    \item[] Answer: \answerNo{}
    \item[] Justification: The benchmark data are public (SWE-bench Verified, SWE-bench Live). The instrument, the sweep driver, the analysis scripts and the archived timelines will be released with the camera-ready version; they are not attached to the anonymous submission.

\item {\bf Experimental setting/details}
    \item[] Question: Does the paper specify all the training and test details (e.g., data splits, hyperparameters, how they were chosen, type of optimizer) necessary to understand the results?
    \item[] Answer: \answerYes{}
    \item[] Justification: No models are trained. Section 4 gives the slices and how they were fixed, the harness and its three recovery policies, the exact model strings, the per-run cost cap, the outcome definitions, and the pre-declared analysis.

\item {\bf Experiment statistical significance}
    \item[] Question: Does the paper report error bars suitably and correctly defined or other appropriate information about the statistical significance of the experiments?
    \item[] Answer: \answerYes{}
    \item[] Justification: Resolve and bearing fractions carry exact (Clopper--Pearson) 95\% intervals; paired comparisons use exact sign tests and exact McNemar tests, the cross-stack comparison a one-sided Fisher exact test with its sensitivity to reassigned runs stated; the tests reported and the absence of multiplicity adjustment are declared (Section 5.6, Appendix B). Run-to-run variability is measured only for plain rollback (Section 5.6), and intervals ignore repository clustering (Section 7).

\item {\bf Experiments compute resources}
    \item[] Question: For each experiment, does the paper provide sufficient information on the computer resources (type of compute workers, memory, time of execution) needed to reproduce the experiments?
    \item[] Answer: \answerYes{}
    \item[] Justification: Section 3 gives replay cost per observation and per sweep; Table 1 gives model spend per arm and substrate; Appendix B states the machine (one 12-vCPU, 22 GB x86 virtual machine running Docker, no GPUs) and the total model spend including pilots.

\item {\bf Code of ethics}
    \item[] Question: Does the research conducted in the paper conform, in every respect, with the NeurIPS Code of Ethics \url{https://neurips.cc/public/EthicsGuidelines}?
    \item[] Answer: \answerYes{}
    \item[] Justification: The work measures software agents on public benchmarks; no human subjects, personal data, or high-risk assets are involved, and the submission is anonymized.

\item {\bf Broader impacts}
    \item[] Question: Does the paper discuss both potential positive societal impacts and negative societal impacts of the work performed?
    \item[] Answer: \answerNo{}
    \item[] Justification: The contribution is an evaluation instrument and measurements of it. The plausible impact is positive (agents that break less existing behaviour); we see no direct path to a negative application, and the paper does not discuss societal impact.

\item {\bf Safeguards}
    \item[] Question: Does the paper describe safeguards that have been put in place for responsible release of data or models that have a high risk for misuse (e.g., pre-trained language models, image generators, or scraped datasets)?
    \item[] Answer: \answerNA{}
    \item[] Justification: No models or scraped datasets are released.

\item {\bf Licenses for existing assets}
    \item[] Question: Are the creators or original owners of assets (e.g., code, data, models), used in the paper, properly credited and are the license and terms of use explicitly mentioned and properly respected?
    \item[] Answer: \answerYes{}
    \item[] Justification: SWE-bench, SWE-bench Verified and SWE-bench Live are cited [1, 2, 3] and used under their MIT licenses and published container images; the models are accessed through their providers' APIs under the providers' terms; every cited artefact appears in the references.

\item {\bf New assets}
    \item[] Question: Are new assets introduced in the paper well documented and is the documentation provided alongside the assets?
    \item[] Answer: \answerNo{}
    \item[] Justification: The instrument and the archived timelines are not released with the anonymous submission; they will be released, documented, with the camera-ready version.

\item {\bf Crowdsourcing and research with human subjects}
    \item[] Question: For crowdsourcing experiments and research with human subjects, does the paper include the full text of instructions given to participants and screenshots, if applicable, as well as details about compensation (if any)?
    \item[] Answer: \answerNA{}
    \item[] Justification: No crowdsourcing or human-subject research.

\item {\bf Institutional review board (IRB) approvals or equivalent for research with human subjects}
    \item[] Question: Does the paper describe potential risks incurred by study participants, whether such risks were disclosed to the subjects, and whether Institutional Review Board (IRB) approvals (or an equivalent approval/review based on the requirements of your country or institution) were obtained?
    \item[] Answer: \answerNA{}
    \item[] Justification: No human subjects.

\item {\bf Declaration of LLM usage}
    \item[] Question: Does the paper describe the usage of LLMs if it is an important, original, or non-standard component of the core methods in this research?
    \item[] Answer: \answerYes{}
    \item[] Justification: LLMs are the agents under measurement: the planner, worker and monitor models are named by their API identifiers in Section 4, and their roles in the harness are described in Sections 3 and 4.

\end{enumerate}
